\documentclass[12pt,a4paper,oneside]{ctexart}
\usepackage{amsmath,amsthm,amssymb,bm,color,framed,graphicx,hyperref,mathrsfs,fancyhdr}

\title{\textbf{算法基础 Homework 1}}
\author{郝思粤 \\ 学号 PB23151779 \\ 中国科学技术大学}
\date{\today}
\linespread{1.5}
\definecolor{shadecolor}{RGB}{241,241,255}

% 页眉页脚
\pagestyle{fancy}
\fancyhf{}
\fancyhead[L]{算法基础 Homework 1}
\fancyhead[R]{PB23151779}
\fancyfoot[C]{\thepage}

% Problem 环境
\newcounter{problemname}
\newenvironment{problem}{
    \begin{shaded}
    \stepcounter{problemname}
    \addcontentsline{toc}{section}{题目\arabic{problemname}}
    \noindent\textbf{题目\arabic{problemname}. }
}{
    \end{shaded}\par
}
\newenvironment{solution}{
    \addcontentsline{toc}{subsection}{解答}
    \noindent\textbf{解答. }
}{\par}
\newenvironment{note}{
    \addcontentsline{toc}{subsection}{注记}
    \noindent\textbf{题目\arabic{problemname}的注记. }
}{\par}


\begin{document}
\maketitle
\tableofcontents
\newpage

\begin{problem}
（ 5 $\times$ 2 = 10 分）判断以下命题的正误。若正确，请给出证明；若错误，请给出反例。

1. 若 \( f(n) = \Theta(g(n)) \)，则有 \( \lg(f(n)) = \Theta(\lg g(n)) \)。

2. 若 \( f(n) = \sum_{i}^{l} b_{i} n^{i} \)，且 \( b_{l} > 0 \)，则有 \( f(n) = \Theta(n^{l}) \).

\end{problem}

\begin{solution} 
\par 1. \textbf{错误（一般不成立）。}
    令对所有 \(n\)，
    \[
    f(n)=2,\qquad g(n)=1.
    \]
    显然 \(f(n)=\Theta(g(n))\)，因为存在常数 \(c_1=1,c_2=2\) 及 \(n_0=1\) 使得对所有 \(n\ge n_0\) 有
    \[
    c_1 g(n) \leq  f(n) \leq c_2 g(n)
    \]
   但
    \[
    \lg f(n)=\lg 2 = 1,\qquad \lg g(n)=\lg 1 = 0.
    \]
    于是不可能存在常数 \(a,b>0\) 使得 \(a\cdot \lg g(n)\le \lg f(n)\le b\cdot \lg g(n)\) 对 足够大的 \(n\) 成立（因为右侧始终为 $0$）。因此 \(\lg f(n)\) 不是 \(\Theta(\lg g(n))\)。

\par 2. \textbf{正确。}
    设
    \[
    f(n)=b_{\ell}n^{\ell}+\sum_{i=k}^{\ell-1} b_i n^i,
    \]
    其中 \(b_{\ell}>0\)。

    \textbf{上界：} 对所有 \(n\ge 1\)，因为当 \(i\le \ell\) 时 \(n^i \le n^{\ell}\)，有
    \[
    f(n) \le \sum_{i=k}^{\ell} |b_i| n^i \le \Big(\sum_{i=k}^{\ell} |b_i|\Big) n^{\ell}.
    \]
    取 \(c_2=\sum_{i=k}^{\ell} |b_i|\)，则 \(f(n)\le c_2 n^{\ell}\)，即 \(f(n)=O(n^{\ell})\)。

    \textbf{下界：} 写作
    \[
    f(n)=b_{\ell}n^{\ell}+\sum_{i=k}^{\ell-1} b_i n^i.
    \]
    令 \(C=\sum_{i=k}^{\ell-1} |b_i|\)，则
    \[
    \Big|\sum_{i=k}^{\ell-1} b_i n^i\Big| \le C n^{\ell-1}.
    \]
    于是
    \[
    f(n) \ge b_{\ell}n^{\ell} - C n^{\ell-1}
    = n^{\ell}\left(b_{\ell}-\tfrac{C}{n}\right).
    \]
    由于 \(b_{\ell}>0\)，取 \(n_0=\max\{1,\lceil 2C/b_{\ell}\rceil\}\)，则当 \(n\ge n_0\) 时有
    \[
    f(n)\ge \frac{b_{\ell}}{2} n^{\ell}.
    \]
    取 \(c_1=b_{\ell}/2\)，即得 \(f(n)\ge c_1 n^{\ell}\)，即 \(f(n)=\Omega(n^{\ell})\)。

    \textbf{结论：} 综上，存在常数 \(c_1,c_2>0\) 与 \(n_0\)，使得对所有 \(n\ge n_0\) 有
    \[
    c_1 n^{\ell} \le f(n) \le c_2 n^{\ell}.
    \]
    因此，\(f(n)=\Theta(n^{\ell})\)。

\end{solution}

\begin{note}
\par 第一小问成立的条件：
\begin{itemize}
    \item \(f(n)\) 和 \(g(n)\) 都是正函数（即对所有 \(n\) 都大于零）。
    \item \(f(n)\) 和 \(g(n)\) 都是单调递增的。
\end{itemize}
\end{note}


\begin{problem}
Q2. (15 + 15 + 10 = 40 分) 函数 \( A: \mathbb{N}^2 \to \mathbb{N} \) 在一些算法的时间复杂度分析中具有重要应用，其定义如下：

\[
A(m, n) = 
\left\{
\begin{array}{ll}
n + 1, & \text{若 } m = 0 \\[1ex]
A(m-1, 1), & \text{若 } m > 0,\, n = 0 \\[1ex]
A(m-1, A(m, n-1)), & \text{若 } m > 0,\, n > 0
\end{array}
\right.
\]


(1) 证明 \( A \) 是良定义的，即对于任意 \( m, n \in \mathbb{N} \)，\( A(m, n) \) 的递归定义总能终止。

(2) 证明 \( A(m, n) \) 关于 \( m, n \) 分别单调递增，即：
\[
\begin{aligned}
A(m + 1, n) &> A(m, n) \\
A(m, n + 1) &> A(m, n)
\end{aligned}
\]

(3) \(\alpha(x)\) 定义为使得 \( A(n, n) \leq x \) 的最大自然数 \( n \)。证明：
\[
\begin{aligned}
\alpha(x) &= \omega(1) \\
\alpha(x) &= O(\lg^* x)
\end{aligned}
\]

其中 \(\lg^* x\) 为迭代对数函数。
\end{problem}

\begin{solution}

\par 1： 证明 $A$ 是良定义的

\textbf{证明（按 $m$ 的完全归纳）}：

 若 $m=0$，按定义
\[
A(0,n)=n+1,\qquad \forall n\ge 0,
\]
此时递归显然终止。故命题对 $m=0$ 成立。

\textbf{归纳假设：} 假设对任意 $k$ 满足 $0\le k\le m$，以及任意 $n\in\mathbb{N}$，$A(k,n)$ 都能终止并给出一个自然数值。
接下来证明命题对 $m+1$ 也成立：即对任意 $n$，$A(m+1,n)$ 终止。

\textbf{对 $n$ 作归纳（在固定的 $m+1$ 下）：}

\begin{itemize}
  \item 当 $n=0$ 时，由定义有
    \[
    A(m+1,0)=A(m,1).
    \]
    右边属于第一个参数不超过 $m$ 的情形，故按归纳假设 $A(m,1)$ 终止，因此 $A(m+1,0)$ 终止。

  \item 归纳步：设对某个 $n\ge0$，已知 $A(m+1,t)$ 对所有 $t\le n$ 都终止。考虑
    \[
    A(m+1,n+1)=A\bigl(m,\,A(m+1,n)\bigr).
    \]
    根据归纳假设（关于 $n$），$A(m+1,n)$ 终止并产生某个自然数值，记作 $y$. 由于 $y\in\mathbb{N}$ 且 $m\le m$，按外层关于 $m$ 的归纳假设（对所有 $k\le m$ 都成立），$A(m,y)$ 也必终止。于是 $A(m+1,n+1)$ 终止。
\end{itemize}

由 $n$ 的归纳原理，固定的 $m+1$ 对所有 $n$ 都成立；即对任意 $n$，$A(m+1,n)$ 终止。

由对 $m$ 的完全归纳，命题对所有 $m$ 成立。故对任意 $m,n\in\mathbb{N}$，$A(m,n)$ 的递归定义均能在有限步内低轨到到m = 0或是n = 0的情况，并可得出确定值，函数 $A$ 定义良好。

\hfill$\square$

\par 2： 证明单调性

\textbf{命题}：对任意 $m,n\in\mathbb{N}$，
\[
A(m+1,n) > A(m,n), \qquad A(m,n+1) > A(m,n).
\]

\medskip
\textbf{证明}：对 $m$ 作完全归纳。

\textbf{当 $m=0$ 时}，
\[
A(0,n)=n+1,\quad A(1,n)=n+2,
\]
显然 $A(1,n)>A(0,n)$，且 $A(0,n+1)=n+2>A(0,n)$。命题成立。

\textbf{归纳假设}：假设对所有 $k\leq m$，命题成立。

\textbf{归纳步}：证明对 $m+1$ 命题成立。

\emph{(i) 关于第二个参数的单调性：}
对 $n$ 归纳。
\[
A(m+1,0)=A(m,1),\quad A(m+1,1)=A(m,A(m+1,0)).
\]
由归纳假设，$A(m,\cdot)$ 随参数严格递增，因而 $A(m+1,1)>A(m+1,0)$。  
若已知 $A(m+1,n+1)>A(m+1,n)$，则
\[
A(m+1,n+2)=A(m,A(m+1,n+1)) > A(m,A(m+1,n))=A(m+1,n+1),
\]
由此得证 $A(m+1,n)$ 关于 $n$ 严格递增。

\emph{(ii) 关于第一个参数的单调性：}
同样对 $n$ 归纳。

当 $n=0$ 时，利用归纳假设中 $A(m,\cdot)$ 的单调性，得到：
  \[
  A(m+2,0)=A(m+1,1)=A(m,A(m,1))>A(m,1)=A(m+1,0),
  \]
  

若已知 $A(m+2,n)>A(m+1,n)$，则
  \[
  A(m+2,n+1)=A(m+1,A(m+2,n)) > A(m,A(m+2,n)) > A(m,A(m+1,n))=A(m+1,n+1),
  \]
  第一不等式由归纳假设（较小的 $m$ 情形）保证，第二不等式由 $A(m,\cdot)$ 的单调性保证。

因此 $A(m+2,n)>A(m+1,n)$ 对任意 $n$ 成立。

综上，命题对所有 $m,n$ 成立。

\hfill$\square$

\par 3： 证明 $\alpha(x)$ 的渐近性质

\textbf{(i) 证明 $\alpha(x)=\omega(1)$}

由定义有，即证，对任意常数 $c>0$，存在 $x_0 > 0$, 使得对所有 $x\ge x_0$，都有$\alpha(x)>c$。
\par 由 $A$ 的递增性，对于$m<n$,有$A(n,n)>A(m,n)>A(m,m)$,不妨设$m=c$,则当$x\ge A(c,c)$时,$\alpha(x)\ge c+1>c$.
\par 取$x_0=A(c,c)$,则对所有$x\ge x_0$,有$\alpha(x)>c$.因此$\alpha(x)=\omega(1)$
\end{solution}

\textbf{(ii) 证明\quad $\alpha(x) = O(\lg^\ast x)$.}

根据定义，$\alpha(x)$ 为满足
\[
A(\alpha(x), \alpha(x)) \leq x
\]
的最大自然数。为了得到 $\alpha(x)$ 的上界，先讨论 $m$较小 时 Ackermann 函数的具体形式：
\[
A(1,n) = n+2, \quad A(2,n) = 2n+3, \quad A(3,n) = 2^{n+3}-3.
\]

对于 $m=4$，递推定义给出：
\[
A(4,n) = A(3, A(4,n-1)).
\]
逐步展开可得
\[
A(4,n) = \mathsf{tower}_2(n+3) - 3,
\]
其中 $\mathsf{tower}_2(k)$ 表示一个以 $2$ 为底、指数高度为 $k$ 的幂塔，即
\[
\mathsf{tower}_2(1) = 2,\quad \mathsf{tower}_2(2) = 2^2,\quad \mathsf{tower}_2(3) = 2^{2^2}, \quad \dots
\]

由 $\alpha(x)$ 的定义，必有
\[
A(\alpha(x), \alpha(x)) \leq x.
\]
当 $\alpha(x) \geq 4$ 时，有
\[
A(\alpha(x), \alpha(x)) \geq A(4,\alpha(x)) = \mathsf{tower}_2(\alpha(x)+3)-3.
\]
因此推出
\[
\mathsf{tower}_2(\alpha(x)+3) - 3 \leq x.
\]

另一方面，迭代对数函数 $\lg^\ast x$ 的定义是：对 $x$ 连续取对数，直至结果不大于 1 所需的次数。它等价于
\[
\mathsf{tower}_2(\lg^\ast x) \geq x.
\]

由不等式
\[
\mathsf{tower}_2(\alpha(x)+3) \leq x \leq \mathsf{tower}_2(\lg^\ast x),
\]
可知
\[
\alpha(x)+3 \leq \lg^\ast x.
\]

于是，
\[
\alpha(x) \leq \lg^\ast x - 3,
\]
说明 $\alpha(x)$ 的增长速度至多与 $\lg^\ast x$ 同阶。因此：
\[
\alpha(x) = O(\lg^\ast x).
\]
\hfill $\blacksquare$



\begin{problem}
Q3.（\( 3 \times 5 = 15 \) 分）假设你有以下列出的五种运行时间的算法。（假设这些是作为输入大小 \( n \) 函数执行的确切操作次数。）假设你有一台计算机，每秒可以执行 \( 10^{10} \) 次操作，你需要在最多一个小时的计算时间内得到结果。对于每种算法，你能在一小时内得到结果的最大输入大小 \( n \) 是多少？

\begin{enumerate}
    \item \( n^4 \)
    \item \( 100n^2 \)
    \item \( n \log n \)
    \item \( 2^n \)
    \item \( 2^{2^n} \)
\end{enumerate}
\end{problem}


\begin{solution}
每小时可执行的总操作次数为
\[
T = 10^{10} \times 3600 = 3.6 \times 10^{13}.
\]

考虑不同时间复杂度函数下的最大可处理输入规模：

\begin{enumerate}
    \item 对于 $f(n) = n^4$，需要满足
    \[
    n^4 \le T < (n+1)^4.
    \]
    两边取四次方根，得到
    \[
    n \le T^{1/4} < n+1 \quad \Longrightarrow \quad n \le (3.6 \times 10^{13})^{1/4} < n+1.
    \]
    因此，最大可处理整数输入规模为
    \[
    n_{\max} = 2449.
    \]
    
    \item 对于 $f(n) = 100 n^2$，不等式为
    \[
    100 n^2 \le T < 100 (n+1)^2 \quad \Longrightarrow \quad n^2 \le 3.6 \times 10^{11} < (n+1)^2.
    \]
    取平方根得到
    \[
    n \le 6 \times 10^5 < n+1,
    \]
    所以最大整数输入规模为
    \[
    n_{\max} = 6 \times 10^5.
    \]
    
    \item 对于 $f(n) = n \log_2 n$，需满足
    \[
    n \log_2 n \le T < (n+1) \log_2 (n+1).
    \]
    由于无法显式求解，采用数值近似法，得到
    \[
    n_{\max} \approx 9.06 \times 10^{11}.
    \]
    即
    \[
    n \le 9.06 \times 10^{11} < n+1.
    \]
    
    \item 对于 $f(n) = 2^n$，不等式为
    \[
    2^n \le T < 2^{n+1} \quad \Longrightarrow \quad n \le \log_2 T < n+1.
    \]
    计算得到
    \[
    n \le 45.1 < n+1 \quad \Longrightarrow \quad n_{\max} = 45.
    \]
    
    \item 对于 $f(n) = 2^{2^n}$，不等式为
    \[
    2^{2^n} \le T < 2^{2^{n+1}}.
    \]
    取两次对数得到
    \[
    n \le \log_2(\log_2 T) < n+1 \quad \Longrightarrow \quad n_{\max} = 5.
    \]
\end{enumerate}


\end{solution}


\begin{problem}(15 + 20 = 35 分)

中秋佳节即将到来，科大糕点厂决定给 \( n \) 个科大幼儿园的小朋友分发科大定制月饼。小朋友们排成一队依次领取，从队头数起第 \( i(1 \leq i \leq n) \) 个小朋友有 \( a_i \) 朵小红花，分到 \( c_i \) 枚月饼。同时，为了奖励小红花多的小朋友，幼儿园园长制定了以下分发规则：

\begin{itemize}
    \item 每个小朋友至少分到 1 枚月饼 (\( c_i \geq 1 \))；
    \item 若小朋友的小红花数多于相邻的小朋友的小红花数，则其分到的月饼数也多于相邻的小朋友 (\( a_i > a_j (1 \leq i, j \leq n, |i - j| = 1) \Rightarrow c_i > c_j \))；
\end{itemize}

\begin{enumerate}
    \item 若 \( n = 5, a_1 = 1, a_2 = 3, a_3 = 3, a_4 = 5, a_5 = 2 \)，请给出一个满足以上规则的月饼分配方案。
    
    \item 现在 \( n, a_i \) 已知，请设计满足以上规则的分发方案（算法），输出每个小朋友分到的月饼数 \( c_i \)，同时使得 \(\sum_{i=1}^n c_i\) 最小。给出该算法的伪代码。（不需要证明算法的最优性）
\end{enumerate}

\end{problem}


\begin{solution}
\par 1.

\textbf{解答：}

直接可以根据规则手工分配,一个满足规则的分配方案为：

\[
\boxed{c = (1,2,2,3,1)}
\]


\par 2.


\textbf{算法思路：}

\begin{enumerate}
    \item 初始化每个小朋友 $c_i = 1$；
    \item 从左往右遍历小朋友,使得其满足规则约束：
    \begin{itemize}
        \item 若 $a_i > a_{i-1}$，则 $c_i = c_{i-1} + 1$
    \end{itemize}
    \item 从右往左遍历小朋友，使得其满足规则约束：
    \begin{itemize}
        \item 若 $a_i > a_{i+1}$，则 $c_i = \max(c_i, c_{i+1}+1)$
    \end{itemize}
    \item 输出 $c = (c_1, c_2, \dots, c_n)$
\end{enumerate}

\textbf{伪代码如下：}

\begin{verbatim}
Input: n, a[1..n]
Output: c[1..n]

1. 初始化 c[i] = 1, i = 1..n

2. 从左往右遍历:
   for i = 2 to n:
       if a[i] > a[i-1]:
           c[i] = c[i-1] + 1

3. 从右往左遍历:
   for i = n-1 downto 1:
       if a[i] > a[i+1]:
           c[i] = max(c[i], c[i+1] + 1)

4. 输出 c[1..n]
\end{verbatim}


\end{solution}




\end{document}
